\chapter*{Notación y glosario de símbolos}
\addcontentsline{toc}{chapter}{Notación y glosario de símbolos}
\markboth{Notación}{Notación}

La siguiente tabla resume la notación utilizada de forma recurrente a lo largo del libro. Notación adicional, específica de un capítulo, se define en el lugar donde se introduce.

\renewcommand{\arraystretch}{1.35}
\begin{center}
\begin{tabular}{@{}ll@{}}
\toprule
\textbf{Símbolo} & \textbf{Significado} \\
\midrule
$\Omega$, $\Gamma$ & Dominio del sólido y su frontera \\
$\Gamma_u$, $\Gamma_t$ & Frontera con condición esencial (Dirichlet) y natural (Neumann) \\
$\bm{u}$, $u_i$ & Vector de desplazamientos y su componente $i$ \\
$\bar{\bm{u}}$ & Desplazamiento prescrito en $\Gamma_u$ \\
$\bm{b}$, $b_i$ & Fuerza de cuerpo por unidad de volumen \\
$\bm{t}$, $\bar{\bm t}$ & Vector de tracción; tracción prescrita en $\Gamma_t$ \\
$\bm{\sigma}$, $\sigma_{ij}$ & Tensor de tensiones (notación tensorial e indicial) \\
$\bm{\varepsilon}$, $\varepsilon_{ij}$ & Tensor de deformaciones infinitesimales \\
$\bm{C}$, $C_{ijkl}$ & Tensor constitutivo elástico (4to orden) \\
$W$ & Potencial de energía de deformación \\
$E$, $\nu$, $G$ & Módulo de Young, razón de Poisson, módulo de corte \\
$\lambda$, $\mu$ & Constantes de Lamé \\
$K$ & Módulo de compresibilidad (bulk modulus) \\
$\delta u$, $\delta W$ & Desplazamiento virtual; trabajo virtual \\
$N_i(z)$ & Función de forma asociada al nodo $i$ \\
$\mathbf{N}$ & Matriz (vector fila) de funciones de forma \\
$\mathbf{B}$ & Matriz de relación deformación-desplazamiento \\
$\mathbf{K}^{(e)}$, $\mathbf{K}$ & Matriz de rigidez elemental y global \\
$\mathbf{F}^{(e)}$, $\mathbf{F}$ & Vector de fuerzas nodales elemental y global \\
$\mathbf{U}$ & Vector de desplazamientos nodales \\
$\mathbf{T}$ & Matriz de transformación de coordenadas (rotación) \\
$r$ & Coordenada natural (paramétrica), $r \in [-1,1]$ \\
$J$ & Jacobiano de la transformación isoparamétrica \\
$\alpha_i$, $r_i$ & Peso y punto de la cuadratura de Gauss \\
$(\cdot)_{,j}$ & Derivada parcial respecto a $x_j$ (notación indicial) \\
$\delta_{ij}$ & Delta de Kronecker \\
\bottomrule
\end{tabular}
\end{center}
